\section*{第 \ref{ch:g6:propdata} 章 --- 正比例与数据}

\begin{solution}{exo:g6:propdata:1}
苹果：商 $\frac52 = 2.5$，$\frac{7.5}{3} = 2.5$，
$\frac{12.5}{5} = 2.5$ --- 商全部相等，成正比例（单价
每千克 $2.50$ 元）。

年龄/身高：$\frac{85}{2} = 42.5$，但 $\frac{103}{4} = 25.75$ --- 商
不相等，不成正比例。
\end{solution}

\begin{solution}{exo:g6:propdata:2}
一个牛角面包：$4.80 \div 4 = 1.20$ 欧元。七个：
$7 \times 1.20 = 8.40$ 欧元。
\end{solution}

\begin{solution}{exo:g6:propdata:3}
比例系数：$18 \div 10 = 1.8$ 欧元每升。于是
$25$ 升的费用为 $25 \times 1.8 = 45$ 欧元；$72$ 欧元可买
$72 \div 1.8 = 40$ 升；$60$ 升的费用为 $60 \times 1.8 = 108$ 欧元。
\end{solution}

\begin{solution}{exo:g6:propdata:4}
对于 $250$ 千米，其耗油量为 $2.5$ 倍于 $100$ 千米的耗油量，因此
需 $2.5 \times 6 = 15$ 升。对于 $350$ 千米，需 $3.5 \times 6 = 21$ 升。若有
$27$ 升：$27 \div 6 = 4.5$ 百千米，即 $450$ 千米。
\end{solution}

\begin{solution}{exo:g6:propdata:5}
$50\,\%$ 的 $84$ 为一半，即 $42$。\quad
$25\,\%$ 的 $200$ 为四分之一，即 $50$。\quad
$10\,\%$ 的 $63$ 为十分之一，即 $6.3$。\quad
$30\,\%$ 的 $70$：$3 \times 7 = 21$。
\end{solution}

\begin{solution}{exo:g6:propdata:6}
$60\,\%$ 的 $450$：$\frac{60}{100} \times 450 = 270$ 名学生在
食堂就餐；$450 - 270 = 180$ 名学生不在食堂用餐。
\end{solution}

\begin{solution}{exo:g6:propdata:7}
借书数少于 $10$ 本的日子有星期二（$8$ 本）和星期四（$6$ 本）。星期三
比星期四多 $17 - 6 = 11$ 本。
\end{solution}

\begin{solution}{exo:g6:propdata:8}
条形图，共五个条形，高度分别为 $14$、$16$、$13$、$17$、$20$
（刻度每 $2$ 或 $5$ 度较合适）。最热的一天：星期五
（$20$ 度）。
\end{solution}

\begin{solution}{exo:g6:propdata:9}
对 $10$ 人，乘以 $\frac{10}{6} = \frac53$：面粉
$450 \times \frac{10}{6} = 750$ 克。鸡蛋：
$3 \times \frac{10}{6} = 5$ --- 恰好是整数。（若不是
整数，则应\emph{向上}取整，以保证够用。）
\end{solution}

\begin{solution}{exo:g6:propdata:10}
\emph{1.} 在 $2$ 小时内行驶 $48$ 千米；半小时内行驶 $12$ 千米；在
$1$ 小时 $30$ 分钟内行驶 $24 + 12 = 36$ 千米。

\emph{2.} $60 \div 24 = 2.5$ 小时，即 $2$ 小时 $30$ 分钟。
\end{solution}

\begin{solution}{exo:g6:propdata:11}
球：降价额为 $20\,\%$ 的 $15 = 3$ 欧元，现价 $12$ 欧元。
球拍：降价 $8$ 欧元，现价 $32$ 欧元。每种情形下，新价格都是
旧价格的 $80\,\%$，成正比例，比例系数为
$0.8$。
\end{solution}

\begin{solution}{exo:g6:propdata:12}
经过 $t$ 小时，粗蜡烛已燃烧了其高度的 $\frac t6$，
即剩余高度为原高度的 $1 - \frac t6$；细蜡烛则为
$1 - \frac t4$。由题意，应有
\[
1 - \frac t4 = \frac12 \left(1 - \frac t6\right)
\quad\text{即}\quad
1 - \frac t4 = \frac12 - \frac t{12}.
\]
于是 $\frac12 = \frac t4 - \frac t{12} = \frac{3t - t}{12} =
\frac{t}{6}$，故 $t = 3$ 小时。检验：经过 $3$ 小时，粗蜡烛
剩余高度为 $1 - \frac36 = \frac12$，细蜡烛为
$1 - \frac34 = \frac14$ --- 恰为粗蜡烛剩余高度的一半。
\end{solution}

\begin{solution}{pb:g6:propdata:1}
\textbf{1.} $100\,000$ 厘米 $= 1\,000$ 米 $= 1$ 千米：地图上的一厘米
代表地面上的一千米。

\textbf{2.} $7.5$ 厘米 $\rightarrow 7.5$ 千米：两村之间；
$12$ 厘米 $\rightarrow 12$ 千米的湖岸。

\textbf{3.} $20$ 千米 $\rightarrow 20$ 厘米（在地图上）。

\textbf{4.} $5$ 千米 $= 500\,000$ 厘米，故比例尺为
$1:500\,000$。徒步地图（$1:100\,000$）更
详细：两点间距离在图上占 $5$ 厘米，而另一张仅占
$1$ 厘米。

\textbf{5.}
\begin{center}
\renewcommand{\arraystretch}{1.2}
\begin{tabular}{l|ccc}
地图距离（厘米） & $1$ & $2$ & $7.5$ \\
\hline
实际距离（千米） & $1$ & $2$ & $7.5$ \\
\end{tabular}
\end{center}
实际距离 $=$ 地图距离 $\times 1$（按这些单位）：
总是\emph{同一}个乘数，这恰好就是
成正比例的定义
（\cref{def:g6:propdata:def}）。这里的比例系数为 $1$
千米每厘米。

\textbf{6.} $4$ m $= 400$ cm，$400 \div 50 = 8$；
$3$ m $= 300$ cm，$300 \div 50 = 6$：图上画出一个
$8$ cm $\times$ $6$ cm 的长方形。

\textbf{7.} 床：$200 \div 50 = 4$ cm，$90 \div 50 = 1.8$ cm。
门洞：$80 \div 50 = 1.6$ cm。

\textbf{8.} 实际面积：$4 \times 3 = 12$ m$^2$。图上面积：
$8 \times 6 = 48$ cm$^2$。换算：
$12$ m$^2 = 12 \times 10\,000 = 120\,000$ cm$^2$，且
\[
120\,000 \div 48 = 2\,500 .
\]
商不是 $50$，而是 $2\,500 = 50 \times 50$。

\textbf{9.} 纸上 1 cm$^2$ 代表实际中一块
长为 $50$ cm、宽为 $50$ cm 的正方形，共含
$50 \times 50 = 2\,500$ cm$^2$。因此长度除以
$50$ 时，面积就除以 $50 \times 50 = 2\,500$：
即面积随比例尺的\emph{平方}而变化。

\textbf{10.} $6 \times 2\,500 = 15\,000$ cm$^2$，且
$15\,000 \div 10\,000 = 1.5$ m$^2$。

\textbf{11.} 图形暗示星期日的销量是\emph{两倍}
（柱子高两倍）。实际上，增量为 $5$ 欧元，
相对 $105$ 欧元，$5 \div 105 \approx 0.048$：约多 $5\,\%$，
而非多 $100\,\%$。经理无视了
检查坐标轴是否从 $0$ 开始的提醒
（见\cref{met:g6:propdata:read}）。

\textbf{12.} 坐标轴从 $0$ 开始时，柱子分别升到 $105$ 和
$110$ 个小格：两根柱子高度几乎相同，
这才是 $5\,\%$ 差异的诚实图像。

\textbf{13.} 上升：增量为 $20$ 枚邮票，起始数量为
$40$：$\frac{20}{40} = 50\,\%$。下降：减少量为 $20$
枚邮票，起始数量为 $60$：$\frac{20}{60} = \frac13 \approx
33\,\%$。同样是 $20$ 枚邮票，起始值不同 --- 百分数
总是\emph{某物}的分数，而
“某物”已经改变。

\textbf{14.} 先减 $20\,\%$：$100 - 20 = 80$ 欧元。
额外的 $10\,\%$ 作用于 $80$：折扣 $8$ 欧元，最终
价格为 $72$ 欧元。总折扣为 $28$ 欧元，相对原价 $100$ 欧元，因此
$28\,\%$ --- 而非 $30\,\%$。第二次折扣作用于
已降低的价格上，所以它对应的欧元数较小：不同量的百分数
不能直接相加。

\textbf{15.} 在该地图上，$1$ cm 代表 $5$ km，因此 $1$
cm$^2$ 代表 $5 \times 5 = 25$ km$^2$（第 9 题的规则）。
森林面积：$3 \times 25 = 75$ km$^2$。长度随
比例尺变化；面积则随其\emph{平方}变化。
\end{solution}
